%---------------------------------------------------------------------------%
%->> Frontmatter
%---------------------------------------------------------------------------%
%-
%-> 生成封面
%-

\maketitle% 生成中文封面
\MAKETITLE% 生成英文封面
%-
%-> 作者声明
%-
\makedeclaration% 生成声明页
%-
%-> 中文摘要
%-
\intobmk\chapter*{摘\quad 要}% 显示在书签但不显示在目录
\setcounter{page}{1}% 开始页码
\pagenumbering{Roman}% 页码符号

同步定位与地图构建（Simultaneous Localization and Mapping, SLAM）是自主移动机器人、无人机及增强现实设备实现环境感知与自主导航的核心技术。视觉 SLAM 后端优化涉及大规模稀疏矩阵构建与线性方程组求解，具有极高的计算复杂度与内存访问压力，传统嵌入式 ARM 处理器难以实时完成。针对这一核心矛盾，本文提出并系统实现了一套基于 FPGA 的软硬协同 SLAM 后端加速方案，将"舒尔消元 + EKF 增量更新"的算法策略映射为可在 Zynq UltraScale+ 平台高效运行的异构硬件架构。

在算法层面，本文以 SchurVINS 框架为基础，利用光束法平差（BA）中 Hessian 矩阵的箭形块稀疏结构，通过舒尔补消元将路标点状态从系统中消去，再以单次 EKF 观测更新替代传统 Levenberg--Marquardt 多轮迭代求解，将后端计算主体归结为雅可比累加与 $3 \times 3$ 矩阵求逆两个规则化计算核，实现了算法复杂度与硬件友好性的统一。

在硬件设计层面，本文完成了三项核心工作：（1）设计了 5 级深度流水线雅可比计算模块，覆盖从世界坐标变换、重投影残差计算到 Hessian 矩阵在线累加的完整链路，并通过切换外参常量实现左右目双相机的单硬件路径复用；（2）提出了基于 8 位稀疏观测位图（SOB）编码的两级动态门控机制，在流水线前端实现状态级旁路与相机级时钟门控，自动跳过零块计算，在典型稀疏率 50\% 场景下将总浮点运算量降低约 45\%；（3）实现了针对 $3 \times 3$ 对称矩阵的并行化 Schur 消元核心，以代数余子式法并行计算逆矩阵并利用对称性将输出写操作减少约 50\%，Hessian 矩阵片上 BRAM 占用仅为 67~KB。

在系统架构层面，本文遵循逻辑密集型任务归 PS、数据密集型规则计算归 PL 的原则，在 PS 端 ARM 处理器上运行 SVO 2.0 视觉前端追踪、滑动窗口状态管理与 EKF 位姿更新，将雅可比生成、Hessian 累加与 Schur 消元完全卸载至 PL 端 FPGA 流水线。三通道 AXI 通信架构与双缓冲乒乓调度策略有效隐藏了片外存储访问延迟。

实验结果表明，在 Zynq UltraScale+ ZCU102 平台上，Schur Kernel 将单次后端优化延迟从 2.84~ms 压缩至 0.76~ms，实现 3.74 倍加速比，系统后端优化频率稳定在 25~Hz 以上。在轨迹精度方面，EuRoC 数据集上 FPGA 方案平均绝对轨迹误差（ATE）为 0.073~m，与全软件基准的 0.075~m 精度等价；TUM-VI 数据集上 FPGA 方案平均 ATE 为 0.152~m，优于软件基准的 0.182~m。在能效方面，单帧能耗仅为 19.1~mJ，BRAM 占用为 67~KB，均为同类方案中最优。户外真实场景的部署测试进一步验证了系统的工程可行性。

\keywords{视觉 SLAM，视觉惯性里程计，FPGA，硬件加速，软硬协同，舒尔补消元，稀疏观测位图}
%-
%-> 英文摘要
%-
\intobmk\chapter*{Abstract}% 显示在书签但不显示在目录

Simultaneous Localization and Mapping (SLAM) is a fundamental technology enabling environment perception and autonomous navigation for mobile robots, Unmanned Aerial Vehicles (UAVs), and Augmented Reality devices. The backend optimization of visual SLAM involves large-scale sparse matrix construction and linear system solving, imposing prohibitive computational complexity and memory access pressure that conventional embedded ARM processors cannot sustain in real time. To address this critical challenge, this thesis proposes and systematically implements a hardware-software co-designed SLAM backend acceleration scheme based on FPGA, mapping the ``Schur complement elimination + EKF incremental update'' algorithmic strategy onto a heterogeneous hardware architecture running efficiently on the Zynq UltraScale+ platform.

At the algorithmic level, building upon the SchurVINS framework, this work exploits the arrowhead block-sparse structure of the Hessian matrix in Bundle Adjustment (BA), employing Schur complement elimination to marginalize landmark states from the system. A single EKF observation update then replaces the traditional multi-iteration Levenberg--Marquardt solver, reducing the backend computation to two regularized kernels: Jacobian outer-product accumulation and $3 \times 3$ matrix inversion, achieving a unified balance between algorithmic complexity and hardware friendliness.

At the hardware design level, three core contributions are made: (1)~A 5-stage deep-pipelined Jacobian computation module is designed, covering the complete chain from world coordinate transformation and reprojection residual calculation to online Hessian matrix accumulation, with stereo camera branch reuse achieved through extrinsic parameter switching on a single hardware datapath. (2)~A two-level dynamic gating mechanism based on 8-bit Sparse Observation Bitmap (SOB) encoding is proposed, implementing state-level bypass and camera-level clock gating at the pipeline front-end to automatically skip null-block computations, reducing total floating-point operations by approximately 45\% under a typical 50\% sparsity scenario. (3)~A parallelized Schur elimination core targeting $3 \times 3$ symmetric matrices is implemented, computing inverse matrices via the cofactor method in parallel and exploiting output symmetry to reduce write operations by approximately 50\%, with the on-chip BRAM footprint for the Hessian matrix compressed to only 67~KB.

At the system architecture level, following the principle of assigning logic-intensive tasks to the PS and data-intensive regular computations to the PL, the ARM processor on the PS side runs the SVO~2.0 visual front-end tracking, sliding window state management, and EKF pose updates, while Jacobian generation, Hessian accumulation, and Schur elimination are fully offloaded to the FPGA pipeline on the PL side. A three-channel AXI communication architecture combined with a double-buffering ping-pong scheduling strategy effectively hides off-chip memory access latency.

Experimental results on the Zynq UltraScale+ ZCU102 platform demonstrate that the Schur Kernel compresses single backend optimization latency from 2.84~ms to 0.76~ms, achieving a 3.74$\times$ speedup with the system backend optimization frequency sustained above 25~Hz. In terms of trajectory accuracy, the FPGA scheme achieves a mean Absolute Trajectory Error (ATE) of 0.073~m on the EuRoC dataset, equivalent to the 0.075~m of the full-software baseline; on the TUM-VI dataset, the FPGA scheme achieves a mean ATE of 0.152~m, outperforming the software baseline of 0.182~m. In terms of energy efficiency, the per-frame energy consumption is only 19.1~mJ and the BRAM footprint is 67~KB, both representing the best among comparable solutions. Outdoor real-world deployment tests further validate the engineering feasibility of the system.

\KEYWORDS{Visual SLAM, Visual-Inertial Odometry, FPGA, Hardware Acceleration, Hardware-Software Co-design, Schur Complement Elimination, Sparse Observation Bitmap}

\pagestyle{enfrontmatterstyle}%
\cleardoublepage\pagestyle{frontmatterstyle}%

%---------------------------------------------------------------------------%
